\documentclass{article}

\usepackage{amsmath,amssymb}
\usepackage{xurl}
\usepackage[colorlinks=true,linkcolor=blue,citecolor=blue,urlcolor=blue]{hyperref}
\setlength{\emergencystretch}{2em}

% Shared commands for notes in docs/paper-gaps/.

\DeclareUnicodeCharacter{2080}{\ifmmode{}_{0}\else\textsubscript{0}\fi}
\DeclareUnicodeCharacter{2081}{\ifmmode{}_{1}\else\textsubscript{1}\fi}
\DeclareUnicodeCharacter{2082}{\ifmmode{}_{2}\else\textsubscript{2}\fi}
\DeclareUnicodeCharacter{2099}{\ifmmode{}_{n}\else\textsubscript{n}\fi}

\newcommand{\Lean}{\textsc{Lean}}
\ExplSyntaxOn
\NewDocumentCommand{\leanid}{m}{
  \ifmmode
    \textsf{\detokenize{#1}}
  \else
    \group_begin:
    \urlstyle{sf}
    \tl_set:Nn \l_tmpa_tl {#1}
    \tl_replace_all:Nnn \l_tmpa_tl {\_} {_}
    \exp_args:NV \path \l_tmpa_tl
    \group_end:
  \fi
}
\ExplSyntaxOff
\providecommand{\tr}{\operatorname{tr}}
\newcommand{\tnleanIssueBase}{https://github.com/LionSR/TNLean/issues/}
\newcommand{\tnleanPrBase}{https://github.com/LionSR/TNLean/pull/}
\newcommand{\tnleanDocBase}{https://sirui-lu.com/TNLean/docs/}
\newcommand{\ghissue}[1]{\href{\tnleanIssueBase#1}{GitHub issue \##1}}
\newcommand{\ghpr}[1]{\href{\tnleanPrBase#1}{PR \##1}}
\newcommand{\doclink}[1]{\href{\tnleanDocBase#1.html}{\path{#1}}}

% Dirac notation and the identity operator, as in blueprint/src/macros/common.tex.
% \providecommand keeps note-local definitions authoritative.
\usepackage{dsfont}
\providecommand{\ket}[1]{|#1\rangle}
\providecommand{\bra}[1]{\langle #1|}
\providecommand{\braket}[2]{\langle #1 | #2 \rangle}
\providecommand{\ketbra}[2]{|#1\rangle\!\langle #2|}
\providecommand{\Id}{\mathds{1}}
\providecommand{\one}{\mathds{1}}
\providecommand{\C}{\mathbb{C}}
\providecommand{\MN}[1]{M_{#1}(\C)}
\providecommand{\MD}{M_D(\C)}

% External referencing for paper sources.  The build helper writes sanitized
% auxiliary files under build/paper-aux/ and excludes duplicated source labels.
\usepackage{xr}

% Import a generated paper auxiliary file with a caller-supplied label prefix.
% The candidate paths support builds from docs/paper-gaps/, the repository root,
% and build/paper-aux/ itself.
\newcommand{\importpaperaux}[2]{%
  \IfFileExists{../../build/paper-aux/#2.aux}{%
    \externaldocument[#1]{../../build/paper-aux/#2}%
  }{%
    \IfFileExists{build/paper-aux/#2.aux}{%
      \externaldocument[#1]{build/paper-aux/#2}%
    }{%
      \IfFileExists{#2.aux}{%
        \externaldocument[#1]{#2}%
      }{}%
    }%
  }%
}

% General external reference macros
\newcommand{\paperref}[2]{\ref{#1-#2}}
\newcommand{\papereqref}[2]{(\ref{#1-#2})}

% CPSV16 source (1606.00608 / MPDO-22-12-17-2.tex) external document and shortcuts
\importpaperaux{cpsv16-}{1606.00608/MPDO-22-12-17-2-tnlean}
\newcommand{\cpsvref}[1]{\paperref{cpsv16}{#1}}
\newcommand{\cpsveqref}[1]{\papereqref{cpsv16}{#1}}

% Verdict marker: \gapnote{<kind>}{<status>}. Kind is the policy
% classification (clarification, local-correction, scope-restriction,
% unfaithful, false-source, open-gap); status is open, wip (elimination
% underway), resolved, or historical. Machine-read by the site index;
% prints a verdict line.
\newcommand{\gapnote}[2]{\par\noindent\textbf{Verdict:} #1 (#2).\par}


\title{Pure zero correlation length and local orthogonality in
Cirac--Perez-Garcia--Schuch--Verstraete 2017}
\author{}
\date{2026-06-10}

\begin{document}
\maketitle
\gapnote{scope-restriction}{open}

\section{Source statement}

Section~\cpsvref{Sec:MPS} of Ref.~\cite{Cirac2016MPDO_arXiv} defines
correlations independent of distance (CID) by requiring that translating one
observable through any allowed separation leaves the two-region expectation
value unchanged.\footnote{Local source:
\path{Papers/1606.00608/MPDO-22-12-17-2.tex:437--445}, Eq.~(\cpsvref{ZCL}).}
For a basis of normal tensors (BNT) $\{A_j\}_j$, the mixed transfer matrices
are
\begin{align}
    \mathbb E_{j,j'}
        &:= \sum_i A_j^i\otimes\overline{A_{j'}^i}.
    \label{eq:cpsv16_pure_zcl_local_orthogonality_scope_mixed_transfer}
\end{align}
The source defines local orthogonality by
\begin{align}
    \mathbb E_{j,j'} &= 0,
    \qquad j\ne j',
    \tag{LO}
    \label{eq:cpsv16_pure_zcl_local_orthogonality_scope_source_lo}
\end{align}
and defines zero correlation length (ZCL) as the conjunction of CID and local
orthogonality.\footnote{Local source:
\path{Papers/1606.00608/MPDO-22-12-17-2.tex:467--479},
Definitions~\cpsvref{DefLO} and~\cpsvref{DefZCL}.}

For a single tensor $A=\{A^i\}_i$, write
\begin{align}
    \mathbb E_A
        &:= \sum_i A^i\otimes\overline{A^i}
    \label{eq:cpsv16_pure_zcl_local_orthogonality_scope_transfer_matrix}
\end{align}
for its transfer matrix. Theorem~\cpsvref{TheoremZCLPure} of
Ref.~\cite{Cirac2016MPDO_arXiv} asserts that every canonical-form tensor
satisfies
\begin{align}
    \mathrm{ZCL}(A)
        &\Longleftrightarrow \mathbb E_A^2=\mathbb E_A.
    \label{eq:cpsv16_pure_zcl_local_orthogonality_scope_source_equivalence}
\end{align}
Its proof first identifies local orthogonality with
Eq.~\ref{eq:cpsv16_pure_zcl_local_orthogonality_scope_source_lo} and then
identifies transfer idempotence with CID under those mixed-sector
equations.\footnote{Local source:
\path{Papers/1606.00608/MPDO-22-12-17-2.tex:1248--1268}.}

\section{Formal boundary}

The formal development adopts a single-tensor local orthogonality
convention:\footnote{Formal predicate: \leanid{MPSTensor.IsLocallyOrthogonal}.}
\begin{align}
    \mathrm{LO}_{\mathrm{single}}(A)
        &\Longleftrightarrow \mathbb E_A^2=\mathbb E_A.
    \label{eq:cpsv16_pure_zcl_local_orthogonality_scope_formal_local_orthogonality}
\end{align}
The corresponding single-tensor ZCL condition is then\footnote{Formal
predicate: \leanid{MPSTensor.IsZCL}.}
\begin{align}
    \mathrm{ZCL}_{\mathrm{single}}(A)
        &\Longleftrightarrow
        (\mathbb E_A^2=\mathbb E_A)\wedge\mathrm{CID}(A).
    \label{eq:cpsv16_pure_zcl_local_orthogonality_scope_formal_zcl}
\end{align}
The formal single-block equivalence therefore proves\footnote{Formal theorem:
\leanid{MPSTensor.zcl_iff_idempotent_transfer}.}
\begin{align}
    (\mathbb E_A^2=\mathbb E_A)\wedge\mathrm{CID}(A)
        &\Longleftrightarrow \mathbb E_A^2=\mathbb E_A.
    \label{eq:cpsv16_pure_zcl_local_orthogonality_scope_single_block_equivalence}
\end{align}
The forward direction follows from
\begin{align}
    \mathbb E_A^2=\mathbb E_A
        &\Longrightarrow \mathbb E_A^n=\mathbb E_A
        \quad\text{for every }n\geq 1,
    \label{eq:cpsv16_pure_zcl_local_orthogonality_scope_positive_power_idempotence}\\
    \mathbb E_A^n=\mathbb E_A
        \quad\text{for every }n\geq 1
        &\Longrightarrow \mathrm{CID}(A).
    \label{eq:cpsv16_pure_zcl_local_orthogonality_scope_idempotence_implies_cid}
\end{align}
The reverse direction of
Eq.~\ref{eq:cpsv16_pure_zcl_local_orthogonality_scope_single_block_equivalence}
holds immediately because its left-hand side already contains transfer
idempotence.

This result does not establish Theorem~\cpsvref{TheoremZCLPure} of
Ref.~\cite{Cirac2016MPDO_arXiv}, whose proof appears at source
lines~1248--1268. That theorem combines unrestricted physical CID with the
BNT-family equations in
Eq.~\ref{eq:cpsv16_pure_zcl_local_orthogonality_scope_source_lo}.

The source definitions now have distinct formal counterparts. Physical
CID\footnote{Formal predicate: \leanid{MPSTensor.IsPhysicalCID}.} quantifies
over two nonempty physical blocks and arbitrary complementary gap lengths,
including zero; it requires the two-region expectation value to remain invariant
when the gaps are redistributed at fixed total length. Physical BNT
ZCL\footnote{Formal predicate: \leanid{MPSTensor.IsPhysicalBNTZCL}.} is the
conjunction of the BNT relation, unrestricted physical CID, and BNT local
orthogonality. While these definitions encode the source CID property and
Definition~\cpsvref{DefZCL} of Ref.~\cite{Cirac2016MPDO_arXiv}, they do not
yield the equivalence stated in Theorem~\cpsvref{TheoremZCLPure}. In particular,
the idempotence equivalence\footnote{Formal theorem:
\leanid{MPSTensor.zcl_iff_idempotent_transfer}.} applies only to the
single-block and virtual-insertion predicates.

\section{BNT predicates and the raw-weight obstruction}

BNT local orthogonality\footnote{Formal predicate:
\leanid{MPSTensor.IsBNTLocallyOrthogonal}.} records
Eq.~\ref{eq:cpsv16_pure_zcl_local_orthogonality_scope_source_lo} for every pair
of distinct sectors. Writing $\mathrm{BNT}(A)$ when the family forms a basis of
normal tensors for $A$, the source predicate is
\begin{align}
    \mathrm{ZCL}_{\mathrm{BNT}}(A)
        &:=
        \mathrm{BNT}(A)
        \wedge \mathrm{CID}_{\mathrm{physical}}(A)
        \wedge
        (\forall j\ne j',\ \mathbb E_{j,j'}=0).
    \label{eq:cpsv16_pure_zcl_local_orthogonality_scope_physical_bnt_zcl}
\end{align}
This physical condition is formalized directly.\footnote{Formal predicate:
\leanid{MPSTensor.IsPhysicalBNTZCL}.} In contrast, the virtual-insertion
predicate\footnote{Formal predicate: \leanid{MPSTensor.IsBNTZCL}.} uses
arbitrary virtual insertions:
\begin{align}
    \mathrm{ZCL}^{\mathrm{virt}}_{\mathrm{BNT}}(A)
        &:=
        \mathrm{BNT}(A)
        \wedge \mathrm{CID}_{\mathrm{virt}}(A)
        \wedge
        (\forall j\ne j',\ \mathbb E_{j,j'}=0).
    \label{eq:cpsv16_pure_zcl_local_orthogonality_scope_virtual_bnt_zcl}
\end{align}
The positive-gap physical predicate is
\begin{align}
    \mathrm{ZCL}^{\mathrm{pos}}_{\mathrm{BNT}}(A)
        &:=
        \mathrm{BNT}(A)
        \wedge \mathrm{CID}^{\mathrm{pos}}_{\mathrm{physical}}(A)
        \wedge
        (\forall j\ne j',\ \mathbb E_{j,j'}=0),
    \label{eq:cpsv16_pure_zcl_local_orthogonality_scope_positive_gap_bnt_zcl}
\end{align}
where $\mathrm{CID}^{\mathrm{pos}}_{\mathrm{physical}}$ requires both
complementary gaps to be strictly positive.\footnote{Formal predicates:
\leanid{MPSTensor.IsPositiveGapBNTZCL} and
\leanid{MPSTensor.IsPositiveGapPhysicalCID}.}
The unrestricted condition implies this positive-gap restriction,\footnote{Formal
implications: \leanid{MPSTensor.isPositiveGapPhysicalCID_of_isPhysicalCID} and
\leanid{MPSTensor.isPositiveGapBNTZCL_of_isPhysicalBNTZCL}.} and the
corresponding BNT implication follows directly.

The unrestricted BNT-level equivalence is not an open proof target because it
fails under the raw-weight canonical-form normalization of
Ref.~\cite{Cirac2016MPDO_arXiv}. Consider the one-letter tensor
\begin{align}
    A^0
        &= \operatorname{diag}\left(1,\frac12\right).
    \label{eq:cpsv16_pure_zcl_local_orthogonality_scope_raw_weight_tensor}
\end{align}
It has a single normal BNT component with raw weights $1$ and $1/2$. Local
orthogonality holds vacuously and the tensor satisfies physical CID, yet its
transfer matrix is not idempotent. The companion note
\path{docs/paper-gaps/cpsv16_pure_zcl_raw_weight_counterexample.tex} details
this counterexample and identifies the invalid inference at source
lines~1248--1251.\footnote{Tracked by \ghissue{2597}.}
The formalized single-block theorem therefore remains a scope restriction: it
proves the one-block idempotence--CID implication without asserting the false
raw-weight BNT equivalence.

\section{Physical correlations and complementary gaps}

The CID definition at Eq.~(\cpsvref{ZCL}) of Ref.~\cite{Cirac2016MPDO_arXiv}
quantifies over physical observables on arbitrary nonempty finite regions.
For a block observable $O$, the insertion map is
\begin{align}
    \mathbb E_O(X)
        &:=
        \sum_{\boldsymbol i,\boldsymbol j}
        \bra{\boldsymbol j}O\ket{\boldsymbol i}
        A^{\boldsymbol i}X(A^{\boldsymbol j})^\dagger,
    \label{eq:cpsv16_pure_zcl_local_orthogonality_scope_observable_insertion}
\end{align}
following the two-observable construction at Eq.~(\cpsvref{Corr}) of
Ref.~\cite{Cirac2016MPDO_arXiv}. For observed blocks $O_1,O_2$ on a periodic
chain with complementary gap lengths $n_1,n_2$, the contraction is the
operator trace on the virtual matrix space:
\begin{align}
    \tr_{M_D(\C)}
    \left(\mathbb E_{O_2}\circ\mathbb E_A^{n_2}
        \circ\mathbb E_{O_1}\circ\mathbb E_A^{n_1}\right).
    \label{eq:cpsv16_pure_zcl_local_orthogonality_scope_periodic_contraction}
\end{align}
This contraction is not the matrix trace obtained by applying the composition
in Eq.~\ref{eq:cpsv16_pure_zcl_local_orthogonality_scope_periodic_contraction}
to the identity matrix; for a general Kraus family, the two traces differ.

The formal positive-gap theorem proves
\begin{align}
    \mathbb E_A^2=\mathbb E_A
        &\Longrightarrow
        \mathrm{CID}^{\mathrm{pos}}_{\mathrm{physical}}(A).
    \label{eq:cpsv16_pure_zcl_local_orthogonality_scope_idempotence_positive_gap_cid}
\end{align}
\footnote{Formal theorem:
\leanid{MPSTensor.isPositiveGapPhysicalCID_of_isTransferIdempotent}.}
This result does not establish the unrestricted CID definition of
Ref.~\cite{Cirac2016MPDO_arXiv}, which includes adjacent regions. For an
adjacent pair, one complementary transfer power in the source expression is
\begin{align}
    \mathbb E_A^0
        &= \operatorname{id}_{M_D(\C)}.
    \label{eq:cpsv16_pure_zcl_local_orthogonality_scope_zero_gap_power}
\end{align}
Transfer idempotence yields
Eq.~\ref{eq:cpsv16_pure_zcl_local_orthogonality_scope_positive_power_idempotence}
only for $n\geq 1$ and does not equate $\mathbb E_A^0$ with $\mathbb E_A$. The
positive-gap theorem is thus a scope restriction rather than the complete
forward implication of Theorem~\cpsvref{TheoremZCLPure} in
Ref.~\cite{Cirac2016MPDO_arXiv}.

The virtual-insertion CID predicate\footnote{Formal predicate:
\leanid{MPSTensor.IsCID}.} quantifies directly over arbitrary virtual
matrices $X,Y$. Its converse implication to idempotence is valid, but it
imposes a condition stronger than physical CID. For the multiplicity-one,
unit-weight direct sum, simultaneous block injectivity supplies the physical
observables required at source lines~1250--1258 of
Ref.~\cite{Cirac2016MPDO_arXiv}. Specifically, the realization theorem selects
a target BNT sector and realizes the rank-one insertion $|R)(l|$ on that sector
via a region of length at most $3D^5$, while giving zero on all other sector
pairs.\footnote{Formal theorem:
\leanid{MPSTensor.IsBNTCanonicalForm.exists_basis_physicalObservableTransfer_rankOne_le_three_totalDim_pow_five}.}
Because this theorem operates on the unweighted basis direct sum rather than
the raw weighted repeated-copy tensor,\footnote{Formal declarations:
\leanid{directSumTensor P.basis} and \leanid{P.toTensor}.} and because the
source equations involve all copy pairs with their respective weight powers,
unrestricted observable realization for the raw tensor remains separate.

A second obstruction occurs at source line~1250, where the proof in
Ref.~\cite{Cirac2016MPDO_arXiv} claims that
\begin{align}
    \mathbb E_1^2\ne\mathbb E_1
        &\Longrightarrow
        \text{there exists a nonzero eigenvalue $\lambda$ with
        $|\lambda|<1$.}
    \label{eq:cpsv16_pure_zcl_local_orthogonality_scope_false_spectral_inference}
\end{align}
This implication fails when the only subleading spectral value is zero and the
generalized zero-eigenspace has a nontrivial nilpotent Jordan part. Thus the
printed proof does not establish the unconditional reverse direction, even for
the multiplicity-one, unit-weight representative. The formalized contradiction
argument takes exactly the nonzero subleading left and right eigenvectors used
at source lines~1250--1262 as an explicit hypothesis, rather than deriving them
from non-idempotence.

\paragraph{Local correction at Theorem~\cpsvref{TheoremZCLPure}.}

The multiplicity-one, unit-weight setting admits an alternate proof that
bypasses
Eq.~\ref{eq:cpsv16_pure_zcl_local_orthogonality_scope_false_spectral_inference}.
Fix a BNT sector $j$, denote its trace-one stationary state by $\rho_j$, and
choose arbitrary matrices $u,v\in M_{D_j}(\C)$. Simultaneous physical-observable
realization yields sector-supported rank-one insertions
\begin{align}
    \mathbb E_{O_1}
        &= |\rho_j)(u|,
    \label{eq:cpsv16_pure_zcl_local_orthogonality_scope_first_rank_one_insertion}\\
    \mathbb E_{O_2}
        &= |v)(1|.
    \label{eq:cpsv16_pure_zcl_local_orthogonality_scope_second_rank_one_insertion}
\end{align}
Sector compression, stationarity of $\rho_j$, and the trace of a rank-one
superoperator evaluate the periodic contraction for all complementary gap
lengths $n_1,n_2$ as
\begin{align}
    \tr_{M_D(\C)}
    \left(\mathbb E_{O_2}\mathbb E_A^{n_2}
        \mathbb E_{O_1}\mathbb E_A^{n_1}\right)
        &= \tr\left(u\mathbb E_j^{n_1}(v)\right).
    \label{eq:cpsv16_pure_zcl_local_orthogonality_scope_rank_one_contraction}
\end{align}
No transpose appears in
Eq.~\ref{eq:cpsv16_pure_zcl_local_orthogonality_scope_rank_one_contraction}, and
the right-hand side does not depend on $n_2$. Comparing positive-gap CID at
$(n_1,n_2)=(1,2)$ and $(2,1)$ gives
\begin{align}
    \tr\left(u\mathbb E_j(v)\right)
        &= \tr\left(u\mathbb E_j^2(v)\right)
        \qquad\text{for all }u,v.
    \label{eq:cpsv16_pure_zcl_local_orthogonality_scope_trace_pairing_identity}
\end{align}
Because the trace pairing is nondegenerate,
\begin{align}
    \mathbb E_j^2 &= \mathbb E_j.
    \label{eq:cpsv16_pure_zcl_local_orthogonality_scope_sector_idempotence}
\end{align}
Local orthogonality eliminates the off-diagonal mixed transfer maps, ensuring
that the multiplicity-one direct sum is idempotent.

This repair requires neither diagonalizability nor semisimplicity at zero, and
it avoids assuming a nonzero subleading eigenvalue. It constitutes a local
correction to source line~1250 rather than a literal formalization of the
printed spectral argument.\footnote{Formal theorem:
\leanid{MPSTensor.IsBNTCanonicalForm.isTransferIdempotent_basisDirectSum_of_isPositiveGapBNTZCL}.}
It does not resolve the raw-weight counterexample or lift the positive-gap
restriction to adjacent sites.

Combining this repair with the fixed-point characterization of commuting parent
ground spaces on all chain lengths yields the representative equivalence for
the basis direct sum.\footnote{Formal theorem:
\leanid{MPSTensor.IsBNTCanonicalForm.isPositiveGapBNTZCL_basisDirectSum_iff_hasNNCPHGroundSpaces}.}
This theorem applies specifically to the multiplicity-one, unit-weight basis
direct sum with positive complementary gaps, leaving repeated copies, arbitrary
raw weights, and adjacent zero gaps outside its scope.

The older virtual-insertion predicate also quantified over positive-definite
fixed points, making it vacuous for tensors lacking one. The periodic-chain
formulation avoids this boundary-state issue entirely. In the unrestricted
raw-weight setting, the source proof also fails when passing from
non-idempotence of the full transfer matrix to that of an individual normal BNT
component, since the full spectrum includes eigenvalues contributed directly by
raw copy weights.

Accordingly, the virtual-insertion predicate\footnote{Formal predicate:
\leanid{MPSTensor.IsBNTZCL}.} serves as the surrogate, while the positive-gap
predicate\footnote{Formal predicate: \leanid{MPSTensor.IsPositiveGapBNTZCL}.}
encodes the physical positive-gap condition. Neither matches
Definition~\cpsvref{DefZCL} of Ref.~\cite{Cirac2016MPDO_arXiv},\footnote{Represented
in the library by \leanid{MPSTensor.IsPhysicalBNTZCL}.} and the raw-weight
counterexample refutes the unrestricted equivalence between that definition and
transfer-matrix idempotence.

\bibliographystyle{alpha}
\bibliography{references}

\end{document}
